% === Impostazione del documento ==========================
\documentclass[11pt, openany]{extbook}

\usepackage[english,italian]{babel}
\usepackage{hyperref}
\usepackage{setspace}
\onehalfspace

%=== Allinameto testo a sinistra =========================
\setlength{\parindent}{0pt}


% === Regolazione dei margini =============================
\addtolength{\oddsidemargin}{30pt}
\addtolength{\evensidemargin}{-30pt}
\usepackage{fancyhdr}
\usepackage{multirow}
\usepackage{multicol} 
\usepackage[section]{placeins}

%=== CODICE  =========================
\usepackage{xcolor}
\definecolor{commentgreen}{RGB}{2,112,10}
\definecolor{eminence}{RGB}{108,48,130}
\definecolor{weborange}{RGB}{255,165,0}
\definecolor{frenchplum}{RGB}{129,20,83}
\definecolor{dkgreen}{rgb}{0,0.6,0}
\definecolor{gray}{rgb}{0.5,0.5,0.5}
\definecolor{mauve}{rgb}{0.58,0,0.82}

\usepackage{listings}
\lstset {
    language= C,
    frame=tb,
    tabsize=4,
    showstringspaces=false,
    numbers=left,
    %upquote=true,
    commentstyle=\color{commentgreen},
    keywordstyle=\color{eminence},
    stringstyle=\color{red},
    basicstyle=\small\ttfamily, % basic font setting
    emph={int,char,double,float,unsigned,void,bool},
    emphstyle={\color{blue}}, % keyword highlighting
    classoffset=1, % starting new class
    otherkeywords={>,<,.,;,-,!,=,~},
    morekeywords={>,<,.,;,-,!,=,~},
    keywordstyle=\color{weborange},
    classoffset=0
}

\lstset{
  language=R,                     % the language of the code
  basicstyle=\footnotesize,       % the size of the fonts that are used for the code
  numbers=left,                   % where to put the line-numbers
  numberstyle=\tiny\color{gray},  % the style that is used for the line-numbers
  stepnumber=1,                   % the step between two line-numbers. If it's 1, each line
                                  % will be numbered
  numbersep=5pt,                  % how far the line-numbers are from the code
  backgroundcolor=\color{white},  % choose the background color. You must add \usepackage{color}
  showspaces=false,               % show spaces adding particular underscores
  showstringspaces=false,         % underline spaces within strings
  showtabs=false,                 % show tabs within strings adding particular underscores
  frame=tb,                   % adds a frame around the code
  rulecolor=\color{black},        % if not set, the frame-color may be changed on line-breaks within not-black text (e.g. commens (green here))
  tabsize=2,                      % sets default tabsize to 2 spaces
  captionpos=b,                   % sets the caption-position to bottom
  breaklines=true,                % sets automatic line breaking
  breakatwhitespace=false,        % sets if automatic breaks should only happen at whitespace
  title=\lstname,                 % show the filename of files included with \lstinputlisting;
                                  % also try caption instead of title
  keywordstyle=\color{blue},      % keyword style
  commentstyle=\color{dkgreen},   % comment style
  stringstyle=\color{mauve},      % string literal style
  escapeinside={\%*}{*)},         % if you want to add a comment within your code
  morekeywords={*,...}            % if you want to add more keywords to the set
} 


% === Integrazione delle figure ===========================
\usepackage{graphicx}
\usepackage{wrapfig}
\graphicspath{{./imgs/}}
\renewcommand{\figurename}{Fig.}
\usepackage{subfigure}
\usepackage[section]{placeins}
\usepackage[nottoc]{tocbibind}


% === Per le tabelle ======================================
\usepackage{tabularx}
\usepackage{multirow}

% === Per la bibliografia multicolonna ====================
\usepackage{etoolbox}
%\usepackage[backend=biber, sorting=none]{biblatex} \addbibresource{bibliografia-tesi.bib}

%\patchcmd{\thebibliography}{\list}{\begin{multicols}{2}\smaller\list}{}{}
%\appto{\endthebibliography}{\end{multicols}}
% === Per il testo di prova =================================
\usepackage{lipsum}

%%%%%%%%%%%%%%%%%%%%%%%%%%%%%%%%%%%%%%%%%%%%%%%%%%%%%%%%%%%
% TESTO DELLA TESI
%%%%%%%%%%%%%%%%%%%%%%%%%%%%%%%%%%%%%%%%%%%%%%%%%%%%%%%%%%%

\begin{document}
\frontmatter
	% === Frontespizio ====================================
	\pagestyle{empty}
	%%%%%%%%%%%%%%%%%%%%%%%%%%%%%%%%%%%%%%%%%%%%%%%%%%%%%%%%%%%
% Frontespizio
%%%%%%%%%%%%%%%%%%%%%%%%%%%%%%%%%%%%%%%%%%%%%%%%%%%%%%%%%%%

\begin{titlepage}
 \begin{center}
     \includegraphics[width=3.5cm]{Learning/09_ReadME_Documentazione/0_frontespizio/imgs/unimol/unimol_color.png}\\
     \vspace{2em}
     {\Large \textsc{Università  degli studi del Molise}}\\
     \vspace{1em}
     %{\Large \textsc{Corso di Laurea Magistrale in}}\\
     \vspace{1em}
      %{\Large \textsc{Sicurezza dei Sistemi Software }}\\
     \vspace{2em}
     \vspace{4em}
     {\LARGE\textbf{
     ReadMe\\ Documentazione \& Convenzioni\\
       \vspace{1em}
     }}
\begin{center}
     \LARGE{Project Management \& Ingegneria del software} 
\end{center}
 \end{center}

\begin{center}
  \vspace{\fill} %TESTO A FINE PAGINA
{\normalsize Anno Accademico 2021/2022}
\end{center}
\end{titlepage}

\clearpage{\pagestyle{empty}\cleardoublepage}

    
	

	
	% === Indice ==========================================
    \tableofcontents
    \listoffigures
	


	% === Capitoli Tesi ===================================
	\pagestyle{plain}
    \mainmatter
	
\chapter{Introduzione}
\label{chap:intro}
\lipsum[0-2]


	% In questo file importiamo tutti i laboratori
\chapter{Lear Regression}

\begin{lstlisting}[language=R]
# Linear Regression ----

set.seed(1)

## Modello
## Y = B0 + B1X + E

n <- 100

B0 <- 500
B1 <- 0.6

X <- runif(n, 1500, 2500)
X <- sort(X)

m <- B0 + B1*X

### Type param: l, b, p
plot(X, m, type="l")

E <- rnorm(n, 0, 50)

Y <- m + E

plot(X, Y, pch=20, col="blue", ylim=c(1000, 2500))
lines(X, m, col="red", lwd=3)





\end{lstlisting}


	% === Bibliografia ====================================
	\newpage
	\bibliographystyle{IEEEtran}
	\bibliography{bibliografia-tesi.bib}

\end{document}

%%%%%%%%%%%%%%%%%%%%%%%%%%%%%%%%%%%%%%%%%%%%%%%%%%%%%%%%%%%
